\chapter{估值：把尾部目标还原为所需条件}

\section{正式估值只覆盖一个合格标的}

正式估值表只纳入雅化集团，因为它同时具备当前价格、官方财务、公司级经营驱动、完整预测、显式外部目标、估值方法与原始PDF。恩捷虽然有显式目标，但公司级量价与客户证据触发估值阻断；中国船舶、江波龙、盛新锂能与天华新能则缺少点目标或未达到翻倍门槛，均不进入正式目标权重。

\begin{exhibitbox}[Exhibit 8：三档价值与正式目标]
\small
\begin{tabularx}{\exhibitboxwidth}{L{2.4cm} R{1.4cm} R{1.4cm} R{1.4cm} R{1.4cm} R{1.7cm} R{1.7cm} X}
\toprule
标的 & 现价 & 悲观 & 基准 & 牛市 & 牛市空间 & 正式目标 & 正式空间 \\
\midrule
雅化集团 & 16.79 & 9.12 & 26.30 & 42.08 & +150.6\% & 25.32 & +50.8\% \\
\bottomrule
\end{tabularx}
\sourcenote{公司股本、2026年7月22日收盘价、原始券商预测、AStock估值模型。}
\end{exhibitbox}

\section{雅化集团估值桥}

雅化集团不再让三档情景共用同一EPS。AStock把民爆利润底仓与锂业务分开，并明确剔除库存收益与套保损益：悲观/基准/牛市分别假设锂盐出货9/12/13万吨、锂业务正常化净贡献0.9/2.05/2.50万元每吨、民爆归母贡献5.0/5.75/6.0亿元，对应归母净利润13.1/30.35/38.5亿元及EPS 1.14/2.63/3.34元。再分别给予8/10/12.6倍PE，得到9.12、26.30、42.08元。

\begin{exhibitbox}[Exhibit 8A：雅化经营驱动至EPS桥]
\scriptsize
\begin{tabularx}{\exhibitboxwidth}{L{1.5cm} R{1.2cm} R{1.4cm} R{1.3cm} R{1.3cm} R{1.3cm} R{1.4cm} R{1.2cm} R{1.2cm} R{1.2cm}}
\toprule
情景 & 锂盐万吨 & 锂价代理万元/吨 & 收入亿元 & 毛利率 & 毛利亿元 & 费税等亿元 & 归母亿元 & EPS元 & PE/价值 \\
\midrule
悲观 & 9 & 12.0 & 120.0 & 20.0\% & 24.0 & 10.9 & 13.1 & 1.14 & 8x/9.12 \\
基准 & 12 & 16.5 & 170.52 & 27.04\% & 46.11 & 15.76 & 30.35 & 2.63 & 10x/26.30 \\
牛市 & 13 & 18.0 & 195.0 & 28.5\% & 55.58 & 17.08 & 38.5 & 3.34 & 12.6x/42.08 \\
\bottomrule
\end{tabularx}
\sourcenote{基准收入、毛利率、出货与民爆利润参考东吴证券；悲观/牛市为AStock情景假设。费税等为毛利至归母的合并桥。}
\qualitynote{吨贡献、锂价代理及牛熊参数不是公司指引；库存收益与套保损益均不计入正常化贡献。}
\end{exhibitbox}

三档概率为30\%/50\%/20\%，经营驱动概率值为24.30元，而不是由同一EPS只改PE得到。审计值计算为：
\[
0.8\times24.302+0.1\times16.79+0.1\times42.00=25.32\text{元}.
\]
正式空间为50.80\%。对基准EPS而言，翻倍门槛33.58元等于12.77倍；这个倍数并不极端，但从当前6.38倍提升至12.77倍仍需要现金流与周期持续性共同证明。读者应使用约25--26元的区间，而非把25.32元视作精确价格。

\section{恩捷股份：零权敏感性，不是正式估值桥}

东吴证券以2027E正常化EPS 5.16元乘20倍PE给出103元目标；6/13/20倍可机械得到30.96/67.08/103.20元。由于出货、客户涨价、实际单平利润、有效利用率和新线良率缺少公司级量化证据，这一组数值全部属于同源卖方敏感性。本报告不为其赋情景概率，不与现价混合，也不发布正式目标。翻倍门槛95.68元等于18.54倍该卖方EPS，说明市场几乎必须完整接受20倍恢复估值。

\section{为什么不用PEG或简单低PE}

2026年利润增速普遍建立在低基数或周期反转上：雅化从低锂价和套保拖累中恢复，恩捷从隔膜价格与利用率低谷恢复，江波龙受存储涨价与库存成本层影响，盛新与天华同样依赖锂价。PEG会把低基数增长误判为结构增长；简单PE则可能把峰值利润误判为永久分母。

本报告优先使用单位经济与现金流校准后的PE。雅化的民爆分部只获得稳定信用，不套锂盐倍数；恩捷的2027EPS只有在出货、单平利润、利用率与现金共同形成公司证据后才有资格进入正式模型；控股平台投资收益和一次性利润不进入成长PE。

\section{市值与股本审计}

雅化集团总股本11.5256亿股，现价对应总市值193.52亿元。恩捷股份在2026年7月21日完成81.34万股限制性股票注销，最新总股本9.8132亿股，现价对应总市值469.46亿元。若沿用Q1旧股本，会轻微高估市值并造成EPS口径漂移。

\begin{sourcequalitybox}[估值治理]
缺少显式点目标者零外部权重；显式目标但期限未披露者，只有通过公司级估值资格门槛才可获得上限10\%的牛市参考权重。雅化符合该条件，恩捷不符合。概率是审慎表达工具，不是统计胜率；正式输出应读作区间，而非精确点位。
\end{sourcequalitybox}

\clearpage
